\chapter{Reproducibility and Traceability}
\label{app:reproducibility-en}

\section{Purpose of the Appendix}

Reproducibility is separated into two levels. The first is \emph{document
reproducibility}: from a specific Git repository state, thesis-facing data contracts
are checked, deterministic LaTeX tables are rendered, and the PDF is built. The
second is \emph{scientific traceability}: dissertation claims are linked either to
tracked scientific summaries or to frozen evidence identities whose integrity was
previously established. These levels must not be conflated: rebuilding the PDF is
not rerunning a scientific experiment.

\section{Canonical Document Build Path}

The primary repository entry point for the Russian dissertation is
\begin{verbatim}
make thesis
\end{verbatim}
and the English rendering is built with
\begin{verbatim}
make thesis-en
\end{verbatim}
Both targets render language-specific tables from the same frozen thesis data before
XeLaTeX: \code{scripts/build_thesis_assets.py} produces the Russian assets and
\code{scripts/build_thesis_assets_en.py} produces the English assets. The underlying
validation checks the claim contract, the thesis-facing QWake C2 summary, SHA-256 bindings to
tracked Stage~1/2/3 evidence, two locally available QWake bindings, and five frozen
identities of sealed QWake C2 materials. Local files are re-hashed; for sealed
materials absent from the ordinary source tree, preserved identities and their format
are checked instead. Only after those checks pass are files under
\code{thesis/generated/} created. Generated files are not committed to Git and do not
constitute independent scientific evidence.

Data contracts can be checked without LaTeX through
\begin{verbatim}
make thesis-check
\end{verbatim}
The canonical CI build additionally records the Git commit and tree, toolchain
versions, PDF signature and size, and SHA-256 of each concrete language artifact.

\section{Evidence Layers Used by the Dissertation}

Two compact machine-readable layers are used:
\begin{itemize}
  \item \code{thesis/data/core_results_verified_summary.json}, a projection of Stage~1/2/3 results. Each tracked source artifact has an expected SHA-256 that is rechecked during the build;
  \item \code{thesis/data/qwake_c2_verified_summary.json}, a compact representation of the sealed QWake C2 result and independently verified aggregates. The build reconciles C08--C11 statuses against the frozen selection result, highest-coverage zero-danger-rule aggregates, and the protocol boundary without rerunning rule evaluation or changing the accounting model. The file is explicitly marked \code{not_new_scientific_evidence}.
\end{itemize}

The complete QWake verification material is intentionally not copied into the ordinary
dissertation source tree. The locally available frozen C1 request and QWake-FP
contract are re-hashed and must match their expected SHA-256 values. For five other C2
materials, frozen artifact, receipt, selected-rule, and common decision-cost-sequence
identities are retained; their source bytes are not claimed to be locally reverified
by this build. This makes the exact sealed result referenced by the text verifiable
without implying that the compact thesis JSON replaces the primary scientific
material.

The build exposes this boundary through separate machine signals:
\code{THESIS_QWAKE_LOCAL_SOURCE_BINDINGS=2/2_PASS} for the two locally re-hashed
files and \code{THESIS_QWAKE_SEALED_IDENTIFIERS_PRESENT=5/5_PASS} for the five
preserved identities. Their joint provenance consistency is reported as
\code{THESIS_PROVENANCE_CONTRACT=PASS}; this does not mean that all seven QWake
materials are present as bytes in the current repository state.

\section{Bound Sources and Checksums}

The following table is generated directly from the thesis data contracts and keeps
the two categories separate: tracked Stage~1/2/3 source artifacts plus two local
QWake sources whose SHA-256 values are recomputed, and five frozen identities of
sealed QWake C2 materials whose source bytes are outside the ordinary local thesis
build surface.

\input{generated/reproducibility_manifest_EN}

\section{QWake Protocol Boundaries}

For QWake C2 the following protocol facts are frozen: no C2 rule was frozen, C3 was
not opened, and neither automatic retry nor authorization reuse is permitted.
Therefore dissertation reproducibility does not include rerunning C2 or transitioning
to C3. Any future experiment on marginal recognizer cost requires a new estimand, a
new protocol, and a new evidence chain; it is not a continuation of the closed C2
invocation.

\section{Minimal Verification of the Resulting PDF}

After a successful build, a concrete artifact can be identified with
\begin{verbatim}
pdfinfo thesis/main_EN.pdf
sha256sum thesis/main_EN.pdf
\end{verbatim}
The SHA-256 identifies that concrete PDF. Byte-identical PDFs across independent build
environments are not claimed by the current process as a scientific result. Content
reproducibility rests on the exact Git source state together with successful data and
provenance-contract checks.

\section{Practical Independent Verification}

An independent reader can verify the document by:
\begin{enumerate}
  \item obtaining the exact Git commit;
  \item running \texttt{make thesis-check} and confirming the claim, arithmetic, QWake protocol-boundary, and core-result binding checks;
  \item running the epistemic-language and bilingual-congruence checks used by CI;
  \item in an environment with XeLaTeX/Biber, running \texttt{make thesis-en};
  \item checking PDF metadata and SHA-256;
  \item for a disputed numerical claim, tracing the generated table to the thesis JSON and then to the referenced source artifact or frozen scientific identity.
\end{enumerate}
This procedure makes the path from prose to evidence explicit while preserving the
boundary between documenting completed research and initiating new scientific
execution.
